% - In the regression setting of the original CSA paper, envelopes are typically convex because the score space is 
% unbounded and the data does not cluster into separated groups\\
% - In our classification setting, the score space is bounded in $[0,1]^K$ and — crucially — true NC vectors for different 
% classes occupy distinct, well-separated regions of this space (example)\\
% - In the classification, 
% % non-convex envelopes are not only permitted but actively beneficial 
% produce tighter prediction sets without sacrificing the coverage guarantee\\
% - The guarantee still holds because it depends only on the exchangeability of the data and the quantile construction, 
% not on the shape of the envelope
% % This is the central geometric motivation for the radial and strip methods developed in section 3.5
% Briefly contrast with collapsed 1D: if the data is genuinely one-dimensional (all dimensions carry the same signal), 
%collapsing to a scalar is sufficient and non-convexity gives no additional benefit

\section{Non-Convex Envelopes in Classification}
\label{sec3.2: non-convex envelopes}
Envelope convexity was identified as one of the main limitations of standard CSA method \cite{ochoarivera2024conformal} 
when applied to classification in Section~\ref{sec2.4.3: csa limitations}. While convex envelopes naturally model continuous 
target spaces in regression, classification score vectors concentrate in disjoint, class-specific clusters across the score 
space. Fitting a single shared convex envelope forces the boundary to bridge across empty intermediate space, resulting in 
uninformatively large prediction sets. To prevent this over extension, we construct dedicated class conditional envelopes 
$E_h$ for each candidate host, this allows the boundaries to closely adapt closely to the non linear structure of true 
infector scores.

% Also very importantly, dropping convexity does not violate statistical validity; the conformal coverage guarantee holds for 
% any measurable envelope $E \subseteq \mathbb{R}_+^K$ under exchangeability \cite{ochoarivera2024conformal}. To ensure that 
% the boundary is both computationally tractable and biologically sensible, the envelope must satisfy two geometric properties:
% \begin{enumerate}
%     \item \textbf{Star shapedness with respect to the origin:} To ensure the scaling function $\tau(\mathbf{s}) = \inf\{t > 0 
%            : \mathbf{s} \in t \cdot E\}$ is continuous and yields a unique scalar threshold.
%     \item \textbf{Coordinate wise monotonicity (convexity toward the axes):} In nonconformity space, the origin $\mathbf{0}$ 
%            represents perfect conformity (maximum infection evidence). If a score vector $\mathbf{x} \in E$ is accepted, any 
%            point $\mathbf{y}$ with $0 \le y_k \le x_k$ for all $k$ exhibits even stronger evidence and must also be accepted.
% \end{enumerate}

Also very importantly, dropping convexity does not violate statistical validity; the conformal coverage guarantee holds for 
any measurable envelope $E \subseteq \mathbb{R}_+^K$ under exchangeability \cite{ochoarivera2024conformal}. 
To ensure that the boundary is both computationally tractable and biologically sensible, the envelope must satisfy two 
geometric properties: star shapedness (Definiton \ref{def: star shaped}) with respect to the origin (ensuring a unique, 
continuous scaling threshold $\tau(s)$) and coordinate wise monotonicity (Definition \ref{def:coordinate-monotone}) which 
ensures convexity toward the axes, so stronger conformity vectors closer to the origin are 
unconditionally included.

\begin{definition}[Star-Shaped Set]
\label{def: star shaped}
A set $E \subseteq \mathbb{R}^K$ is said to be \emph{star-shaped with respect to the origin} $\textbf{0} \in \mathbb{R}^K$ if 
for every point $x \in E$ and every scalar $\lambda \in [0, 1]$, the line segment connecting the origin to $x$ lies entirely 
within $E$:
$$\lambda x \in E \quad \forall x \in E, \; \forall \lambda \in [0, 1].$$
\end{definition}


\begin{definition}[Coordinate Wise Monotone Set]
\label{def:coordinate-monotone}
A set $E \subseteq \mathbb{R}_+^K$ is coordinate wise monotone if for all $\mathbf{x} \in E$ and $\mathbf{y} \in 
\mathbb{R}_+^K$:
\begin{equation*}
\mathbf{0} \le \mathbf{y} \le \mathbf{x} \implies \mathbf{y} \in E,
\end{equation*}
where $\mathbf{y} \le \mathbf{x}$ denotes $y_k \le x_k$ for all coordinates $k \in \{1, \dots, K\}$.
\end{definition}




% While convex envelopes naturally model 
% continuous target spaces in regression, classification score vectors concentrate in disjoint, class-specific clusters across 
% the score space. Therefore, it leads to the envelope being stretched over regions of empty spaces, making the prediction sets 
% unnecessarily large. Here, we will describe the alternate method that we have used in this report.

% Formally, let $\mathcal{C}_h \subseteq \mathbb{R}^K_+$ denote the support of the transformed score vector 
% $\phi_h(\mathbf{s})$ for true label $H = h \in \mathcal{H}$. In multi-class classification, these supports are strictly 
% pairwise separated:
% \begin{equation}
% \label{eq: separated supports}
% \mathcal{C}_h \cap \mathcal{C}_{h'} = \emptyset, \qquad \operatorname{dist}(\mathcal{C}_h, \mathcal{C}_{h'}) > 0, \qquad 
% \forall h \neq h' \in \mathcal{H}.\end{equation}

% A single convex envelope $E_{\text{conv}} = \operatorname{conv}\left(\bigcup_{h \in \mathcal{H}} \mathcal{C}_h \right)$ 
% must contain every line segment connecting points in $\mathcal{C}_h$ and $\mathcal{C}_{h'}$. These line segments inevitably 
% traverse unpopulated space. Let us define the empty gap as $\text{gap} := E_{\text{conv}} \setminus \bigcup_{h \in \mathcal{H}} 
% \mathcal{C}_h$, the Lebesgue volume decomposes as:
% \begin{equation}
% \label{eq: volume inequality}
% \operatorname{Vol}(E_{\text{conv}}) = \sum_{h \in \mathcal{H}} \operatorname{Vol}(\mathcal{C}_h) + \operatorname{Vol}(\text{gap}), 
% \qquad \text{where } \operatorname{Vol}(\text{gap}) > 0.
% \end{equation}


% Test points falling in this gap are included in the prediction set by construction, $|C(\mathbf{s})|$, and they make 
% them larger without improving coverage. This inefficiency scales with the number of classes $N = |\mathcal{H}|$, making shared 
% convex hulls particularly suboptimal for hostlevel classification with multiple classes ($N > 2$). 
% This mirrors our phage data, where the score vectors cluster tightly around distinct class centroids, this is consistent with 
% the score separation established in Section~\ref{sec2.2.3: limitations}.

% In classification, score vectors for different classes occupy separated regions in the score space. If we attempt to fit a 
% single shared convex envelope, we would force the boundary to bridge across empty intermediate space, which would lead to 
% uninformatively large prediction sets. To prevent this, we fit dedicated class conditional envelopes $E_h$ for each 
% candidate host. This allows the boundary to closely adapt to the non linear structure of true infector scores without the 
% convex over extension.

% Also very importantly, dropping convexity does not violate the statistical validity: the marginal coverage guarantee holds for 
% any measurable envelope $E \subseteq \mathbb{R}^K_+$ by exchangeability \cite{ochoarivera2024conformal}.
% To ensure the scaling function $\tau(\mathbf{s}) = \inf \{ t > 0 : \mathbf{s} \in t \cdot E \}$ yields a single, well defined 
% scalar threshold, $E$ needs only to be star shaped (Definition \ref{def: star shaped}) with respect to the origin:
% \begin{equation*}
% \mathbf{x} \in E \implies \lambda \mathbf{x} \in E, \qquad \forall \lambda \in [0,1].
% \end{equation*}



% Star-shapedness guarantees that $\{t > 0 : \mathbf{s} \in t \cdot E\}$ forms a continuous interval, preserving simple 
% single-pass calibration while allowing the envelope to tightly hug non-convex cluster geometries.

